\setlength{\footskip}{8mm}

\chapter{LITERATURE REVIEW}

This chapter reviews the main theoretical and technical foundations related to this thesis. It first discusses therapeutic communication, client resistance, readiness for change, Motivational Interviewing, and Cognitive Behavioral Therapy as the clinical basis for MI-CBT response selection. It then reviews prior LLM-based counseling systems, including approaches that use memory, stage awareness, strategy selection, and therapeutic evaluation. Finally, the chapter connects these works to search-based and reward-guided methods, highlighting the gap that motivates AIM-RGL: the need for a system that can compare possible near-future client reactions before selecting a therapeutic response.

\section{Clinical Foundations of MI-CBT Therapeutic Dialogue}

Therapeutic dialogue differs from ordinary conversation because the therapist's response should support understanding, collaboration, and meaningful change. A counseling response is not evaluated only by relevance or fluency, but also by whether it strengthens the therapeutic relationship, respects client autonomy, and fits the client's current psychological state. This is especially important when the client shows resistance, ambivalence, or low readiness for change, because poorly timed responses may reduce engagement or increase defensiveness \cite{patterson1985therapist,patterson1994functional,beutler2002resistance}. Figure~\ref{fig:therapeutic-dialogue-loop} illustrates how client state, response timing, therapeutic outcomes, and the next client state are connected.

\begin{figure}[htbp]
    \centering
    \begin{tikzpicture}[
        node distance=0.8cm,
        every node/.style={font=\footnotesize, align=center},
        mainbox/.style={
            rectangle,
            rounded corners,
            draw=black,
            minimum width=4.8cm,
            minimum height=0.8cm,
            text width=4.8cm,
            inner sep=4pt
        },
        warnbox/.style={
            rectangle,
            rounded corners,
            draw=black,
            dashed,
            minimum width=4.1cm,
            minimum height=0.8cm,
            text width=4.1cm,
            inner sep=4pt
        },
        arrow/.style={-{Latex[length=3mm]}, thick}
    ]

    \node[mainbox] (state) {Client state\\Resistance, ambivalence, and readiness};
    \node[mainbox, below=of state] (timing) {Response timing and style\\MI-style support or CBT-style structure};
    \node[mainbox, below=of timing] (response) {Therapeutic response};
    \node[mainbox, below=of response] (outcome) {Interaction outcome\\Alliance, engagement, change talk, and sustain talk};
    \node[mainbox, below=of outcome] (nextstate) {Next client state};

    \node[warnbox, right=1.8cm of timing] (risk) {Poorly timed response\\Defensiveness, disengagement, or weakened alliance};

    \draw[arrow] (state) -- (timing);
    \draw[arrow] (timing) -- (response);
    \draw[arrow] (response) -- (outcome);
    \draw[arrow] (outcome) -- (nextstate);
    \draw[arrow] (nextstate.west) .. controls +(-2.2,0) and +(-2.2,0) .. (state.west);
    \draw[arrow] (timing.east) -- (risk.west);

    \end{tikzpicture}
    \caption{Conceptual relationship between client state, response timing, therapeutic outcomes, and the next client state in therapeutic dialogue.}
    \label{fig:therapeutic-dialogue-loop}
\end{figure}

\subsection{Therapeutic Alliance and Client-Centered Communication}

Therapeutic alliance refers to the collaborative relationship between the therapist and the client. Bordin described the working alliance as involving agreement on goals, agreement on therapeutic tasks, and the emotional bond between therapist and client \cite{bordin1979alliance}. Rogers also emphasized empathy, genuineness, and unconditional positive regard as core conditions for therapeutic personality change \cite{rogers1957necessary}. These foundations suggest that therapeutic dialogue depends not only on the intervention itself, but also on how the intervention is communicated.

Client-centered communication is important in AI-assisted counseling because an LLM may generate a response that sounds supportive but still feels directive, assumptive, or poorly matched to the client's readiness. For example, a response may appear warm while still pushing the client toward action too quickly. In MI-CBT dialogue, this is important because structured CBT techniques may require sufficient trust and collaboration before they are introduced. Therefore, therapeutic alliance is treated as a key quality dimension in this thesis when evaluating whether a response supports engagement and readiness for further therapeutic work.

\subsection{Client Resistance, Ambivalence, and Readiness for Change}

Client resistance refers to reluctance, opposition, or disengagement during the therapeutic process. It may appear openly through arguing or rejecting suggestions, or more subtly through short answers, avoidance, topic changes, or passive agreement \cite{patterson1985therapist,patterson1994functional}. Resistance has also been discussed as a process that may be influenced by therapist behavior, rather than only as a fixed client trait \cite{beutler2002resistance}. In MI, resistance is often understood as a signal that the therapeutic approach may not fit the client's current readiness or concerns \cite{miller2013motivational}.

Ambivalence occurs when clients hold both reasons for and reasons against change. Motivational Interviewing addresses this tension by helping clients explore ambivalence and strengthen their own reasons for change \cite{miller2013motivational,bischof2021motivational}. Therapist responses can also influence whether clients produce more change talk or sustain talk during MI sessions \cite{moyers2006therapist}. In this thesis, ambivalence is important because it creates moments where the system must decide whether the client needs further motivational support or is ready for structured CBT intervention.

Readiness for change explains why the same therapeutic response may help one client but hinder another. The Transtheoretical Model describes change as a staged process that includes precontemplation, contemplation, preparation, action, and maintenance \cite{prochaska1997transtheoretical,raihan2023stages}. A client in precontemplation may not yet view the issue as a problem, while a client in preparation may be more open to concrete strategies or behavioral steps. Therefore, therapeutic response selection should be sensitive to the client's current readiness rather than applying the same intervention style across all moments.

\subsection{Motivational Interviewing, Cognitive Behavioral Therapy, and Their Integration}

Motivational Interviewing is a counseling approach that supports motivation, autonomy, and readiness for change. It emphasizes collaboration, evocation, acceptance, and compassion, and commonly uses open-ended questions, affirmations, reflections, and summaries \cite{miller2013motivational,bischof2021motivational}. MI is especially useful when the client is ambivalent or resistant because it avoids direct confrontation and helps the client explore their own reasons for change. In this thesis, MI provides the foundation for response styles that reduce defensiveness and support readiness.

Cognitive Behavioral Therapy provides a more structured approach for identifying and modifying unhelpful thoughts, beliefs, and behaviors. CBT includes techniques such as identifying automatic thoughts, examining evidence, generating alternative interpretations, and planning behavioral change \cite{beck2020cognitive,chand2023cognitive}. These techniques are useful when the client is ready to reflect and take concrete steps, but they may feel directive if introduced too early. Therefore, CBT is most suitable when the client shows sufficient openness and the therapeutic relationship is stable.

The integration of MI and CBT has been explored in conventional mental health treatment, especially for clients who experience ambivalence or resistance during CBT. Prior work suggests that MI can support treatment engagement and prepare clients for more structured CBT work \cite{arkowitz2004integrating,westra2009adding,westra2016integrating}. Reviews also indicate that MI may be useful as an adjunct to CBT, although the timing and manner of integration remain important \cite{randall2017motivational,marker2018efficacy}. Table~\ref{tab:mi-cbt-comparison} summarizes the main differences between MI and CBT and their relevance to AIM-RGL.

\begin{table}[htbp]
    \caption{Comparison between MI and CBT.}
    \label{tab:mi-cbt-comparison}
    \small
    \begin{tabularx}{\textwidth}{p{0.15\textwidth}X X X}
        \toprule
        \textbf{Aspect} & \textbf{MI} & \textbf{CBT} & \textbf{Relevance to AIM-RGL} \\
        \midrule
        Purpose 
        & Builds motivation, autonomy, and readiness. 
        & Supports thought and behavior change. 
        & Helps choose between motivational support and structured intervention. \\

        Stance 
        & Collaborative and autonomy-supportive. 
        & Structured and skill-focused. 
        & Response style should match resistance and readiness. \\

        Best suited for 
        & Ambivalence, resistance, or uncertainty. 
        & Readiness for reflection, belief testing, or action. 
        & Supports stage-sensitive response selection. \\

        Risk if mistimed 
        & May slow structured progress if overused. 
        & May feel directive if introduced too early. 
        & Motivates future-aware MI-CBT timing. \\
        \bottomrule
    \end{tabularx}
\end{table}

Overall, these clinical foundations show that effective MI-CBT dialogue depends on resistance, ambivalence, readiness, and therapeutic alliance. MI and CBT offer complementary response styles, but their usefulness depends heavily on timing. This motivates AIM-RGL's focus on selecting responses not only by current-turn suitability, but also by their likely near-future effect on the simulated client's therapeutic state.

\section{LLM-Based Counseling Systems and Therapeutic Evaluation}

LLM-based counseling systems have recently been explored as tools for generating supportive, empathetic, and therapy-informed dialogue. However, counseling dialogue requires more than fluent response generation because the system must also consider therapeutic safety, client readiness, alliance, and the appropriateness of the intervention strategy \cite{hua2025scopingreview,nguyen2025llmcounseling}. For this thesis, prior systems are reviewed in terms of how they support CBT, MI, hybrid MI-CBT interaction, client simulation, and therapeutic evaluation.

\subsection{LLM-Based CBT, MI, and Hybrid Counseling Agents}

Existing LLM-based counseling systems have explored different ways of supporting therapy-informed dialogue, including CBT-style response generation, MI-based behavior-change support, counseling dialogue datasets, client simulation, and hybrid MI-CBT interaction. CBT-oriented works usually focus on structured therapeutic skills such as identifying thoughts, cognitive reframing, guided discovery, or generating multi-turn CBT-style counseling conversations. For example, CACTUS constructs multi-turn counseling dialogues grounded in CBT and client personas, while HealMe focuses on cognitive reframing through psychotherapy-informed dialogue \cite{lee2024cactus,xiao2024healme}. Psych8k and ChatCounselor further show the value of domain-specific counseling conversations for training mental health support models, as they use real counseling conversations to improve counseling-related response quality \cite{liu2023chatcounselor}.

MI-based and behavior-change systems focus more on motivation, autonomy, readiness, and the client's movement toward change. CAMI is an MI-grounded counselor agent that uses state inference, topic exploration, action selection, and response generation to support behavior-change counseling \cite{yang2025cami}. CCSMI extends this direction by focusing on consistent client simulation for MI-based counseling, where the simulated client's mental state, motivation, beliefs, plans, and receptivity are tracked across the conversation \cite{yang2025consistent}. These works are relevant to this thesis because they show how LLM-based counseling systems can represent client state and readiness during multi-turn interaction.

Other studies focus on constructing or reconstructing counseling dialogues for training and evaluation. SMILE converts single-turn mental health support data into multi-turn conversations, while CPsyCoun reconstructs Chinese psychological counseling dialogues from counseling reports and proposes an evaluation framework for multi-turn counseling \cite{qiu2024smile,zhang2024cpsycoun}. PATIENT-\(\Psi\) is also relevant because it uses CBT-based cognitive models to simulate patients for training mental health professionals, making it closely related to cognitive conceptualization and simulated client design \cite{wang2024patientpsi}. MAGNeT further extends synthetic counseling data generation by using coordinated specialist agents to generate multi-turn mental health counseling sessions \cite{mandal2025magnet}.

More recent work has also explored adaptive and hybrid counseling agents. Persona Drift studies stage-aware adaptation in CBT chatbots, while ChatThero explores multi-session LLM-supported therapeutic support for behavior change \cite{yun2025personadrift,wang2025chatthero}. FLASH is especially close to this thesis because it combines MI and CBT and uses Flashbulb Memory to guide strategy selection for managing client resistance \cite{anonymous2026flash}. However, these systems mainly rely on current-state signals, stage adaptation, synthetic dialogue construction, or memory from previous interactions. Less attention has been given to explicitly comparing possible near-future client reactions before selecting a therapeutic response, which motivates the future-aware response selection direction of this thesis.

Table~\ref{tab:mi-cbt-llm-systems} summarizes the LLM-based systems most closely related to MI-CBT coordination and adaptive therapeutic strategy selection.

\begin{table}[htbp]
    \caption{LLM-based systems related to MI-CBT coordination and adaptive therapeutic response selection.}
    \label{tab:mi-cbt-llm-systems}
    \small
    \begin{tabularx}{\textwidth}{p{0.18\textwidth}p{0.22\textwidth}X X}
        \toprule
        \textbf{System} & \textbf{Therapeutic Focus} & \textbf{Main Mechanism} & \textbf{Limitation / Gap} \\
        \midrule

        ChatThero \cite{wang2025chatthero}
        & Behavior change with MI and CBT strategies
        & Uses a multi-session, stressor-aware, memory-persistent language agent for addiction recovery support.
        & Uses adaptive therapeutic strategies, but is less focused on explicitly comparing short-horizon future client reactions before response selection. \\

        Persona Drift \cite{yun2025personadrift}
        & Stage-aware CBT adaptation
        & Adjusts CBT chatbot flow when persona or stage drift is detected.
        & Relevant to stage-aware adaptation, but it is not a full MI-CBT coordination framework and does not directly compare future client trajectories. \\

        FLASH \cite{anonymous2026flash}
        & Hybrid MI-CBT for resistance management
        & Uses readiness estimation and Flashbulb Memory to dynamically guide switching between MI and CBT strategies.
        & Adapts mainly from observed client history and accumulated memory, rather than explicitly comparing multiple simulated near-future client reactions before response selection. \\

        AIM-RGL
        & Hybrid MI-CBT with reward-guided lookahead
        & Compares candidate MI-CBT responses and simulated near-future client reactions before selecting a response.
        & Proposed framework evaluated in simulated counseling dialogue. \\

        \bottomrule
    \end{tabularx}
\end{table}

\subsection{Client Simulation and Cognitive Conceptualization}

Client simulation is useful for evaluating LLM-based counseling systems before involving real users, especially in sensitive mental health-related settings. In this thesis, simulated clients provide controlled interaction scenarios involving resistance, ambivalence, and different levels of readiness. Cognitive conceptualization can support this simulation by representing the client's situation, beliefs, emotions, automatic thoughts, coping patterns, and behavioral responses. This allows the simulated client to respond more consistently across turns instead of producing isolated or random replies.

Figure~\ref{fig:client-simulation-loop} illustrates the general simulation loop used to evaluate therapeutic dialogue. The therapist-like system generates a response based on the conversation history, the simulated client produces a reaction based on the client profile and current state, and the resulting dialogue is evaluated using therapeutic metrics.

\begin{figure}[htbp]
    \centering
    \begin{tikzpicture}[
        node distance=1.2cm,
        every node/.style={font=\small, align=center},
        box/.style={
            rectangle,
            rounded corners,
            draw=black,
            minimum width=5.3cm,
            minimum height=0.9cm,
            text width=5.3cm,
            inner sep=6pt
        },
        arrow/.style={-{Latex[length=3mm]}, thick}
    ]

    \node[box] (profile) {Client profile and cognitive conceptualization};
    \node[box, below=of profile] (history) {Conversation history and current client state};
    \node[box, below=of history] (therapist) {Therapist-like system response};
    \node[box, below=of therapist] (client) {Simulated client response};
    \node[box, below=of client] (eval) {Therapeutic evaluation};

    \draw[arrow] (profile) -- (history);
    \draw[arrow] (history) -- (therapist);
    \draw[arrow] (therapist) -- (client);
    \draw[arrow] (client) -- (eval);
    \draw[arrow] (client.west) .. controls +(-2.4,0) and +(-2.4,0) .. (history.west);

    \end{tikzpicture}
    \caption{General client simulation loop for therapeutic dialogue evaluation.}
    \label{fig:client-simulation-loop}
\end{figure}

\subsection{Evaluation Metrics for Therapeutic Dialogue}

Therapeutic dialogue evaluation should consider counseling-specific qualities rather than only general language quality. In this thesis, evaluation focuses on whether responses support CBT competence, MI consistency, therapeutic alliance, and movement toward change. CTRS is used to assess CBT-related therapist competence, MITI is used to assess MI fidelity, WAI is used to assess therapeutic alliance, and change talk and sustain talk are used to examine whether the simulated client's language moves toward or away from change \cite{young1980cognitive,moyers2016miti,horvath1989wai,miller2013motivational,miller2009toward}. These metrics are relevant because AIM-RGL aims to improve not only response fluency, but also therapeutic timing, collaboration, and simulated client readiness.

\begin{table}[htbp]
    \caption{Evaluation metrics for therapeutic dialogue.}
    \label{tab:therapeutic-evaluation-metrics}
    \small
    \begin{tabularx}{\textwidth}{p{0.18\textwidth}X X}
        \toprule
        \textbf{Metric} & \textbf{Source / Basis} & \textbf{Relevance to This Thesis} \\
        \midrule

        CTRS 
        & Cognitive Therapy Rating Scale for evaluating CBT therapist competence \cite{young1980cognitive,beckinstitute2021ctrs}. 
        & Evaluates CBT-style response quality, including structure, guided discovery, and cognitive intervention. \\

        MITI 
        & Motivational Interviewing Treatment Integrity code for assessing MI fidelity \cite{moyers2016miti,moyers2015miti421}. 
        & Evaluates whether responses follow MI principles such as empathy, autonomy support, and MI-consistent behavior. \\

        WAI / Alliance 
        & Working Alliance Inventory for measuring therapeutic alliance \cite{horvath1989wai,hatcher2006wai}. 
        & Evaluates trust, collaboration, goal agreement, and task or approach agreement. \\

        Change Talk 
        & MI construct describing client language favoring change \cite{miller2013motivational,miller2009toward}. 
        & Indicates movement toward readiness, motivation, or action. \\

        Sustain Talk 
        & MI construct describing client language favoring the current state \cite{miller2013motivational,miller2009toward}. 
        & Indicates resistance, ambivalence, or lower readiness for change. \\

        \bottomrule
    \end{tabularx}
\end{table}

Because therapeutic alliance is central to this thesis, the alliance dimension is further represented through three components: goal agreement, approach or task agreement, and affective bond. These dimensions reflect whether the client and counselor share an understanding of therapy goals, agree on how therapeutic work should proceed, and maintain trust and emotional connection \cite{bordin1979alliance,horvath1989wai,hatcher2006wai}. Table~\ref{tab:wai-alliance-items} summarizes the alliance-related questionnaire items used to evaluate these dimensions.

\begin{table}[htbp]
    \caption{Therapeutic alliance dimensions and questionnaire items.}
    \label{tab:wai-alliance-items}
    \small
    \renewcommand{\arraystretch}{1.12}
    \begin{tabularx}{\textwidth}{p{0.24\textwidth}X p{0.08\textwidth}}
        \toprule
        \textbf{Dimension} & \textbf{Question} & \textbf{No.} \\
        \midrule

        \multirow{4}{0.24\textwidth}{Goal}
        & There is mutual understanding about what participants are trying to accomplish in therapy. 
        & Q1 \\

        
        & The client and counselor are working on mutually agreed upon goals. 
        & Q2 \\

        
        & The client and counselor have the same ideas about the client's real problems. 
        & Q3 \\

        
        & The client and counselor have established a good understanding of the changes that would be good for the client. 
        & Q4 \\

        \midrule

        \multirow{4}{0.24\textwidth}{Approach}
        & There is agreement about the steps taken to help improve the client's situation. 
        & Q5 \\

        
        & There is agreement about the usefulness of the current activity in therapy. 
        & Q6 \\

        
        & There is agreement on what is important for the client to work on. 
        & Q7 \\

        
        & The client believes that the way they are working with the problem is correct. 
        & Q8 \\

        \midrule

        \multirow{4}{0.24\textwidth}{Affective Bond}
        & There is a mutual liking between the client and counselor. 
        & Q9 \\

        
        & The client feels confident in the counselor's ability to help the client. 
        & Q10 \\

        
        & The client feels that the counselor appreciates him/her as a person. 
        & Q11 \\

        
        & There is mutual trust between the client and counselor. 
        & Q12 \\

        \bottomrule
    \end{tabularx}
\end{table}

In addition to rating-based metrics, change talk and sustain talk can be represented as proportional indicators of the simulated client's movement toward or away from change. The proportions are defined as follows:

\begin{equation}
\text{Change Talk Ratio} =
\frac{\text{Number of change talk utterances}}
{\text{Total number of client utterances}}
\label{eq:change-talk-ratio}
\end{equation}

\begin{equation}
\text{Sustain Talk Ratio} =
\frac{\text{Number of sustain talk utterances}}
{\text{Total number of client utterances}}
\label{eq:sustain-talk-ratio}
\end{equation}

Together, these evaluation dimensions allow the study to assess whether AIM-RGL improves therapeutic response quality in terms of CBT skill, MI consistency, therapeutic alliance, readiness, and simulated client movement across the dialogue.

\section{Future-Aware Response Selection and Research Gap}

Future-aware response selection refers to choosing an action by considering not only its immediate suitability, but also its possible downstream effects. This idea is common in artificial intelligence, especially in sequential decision-making problems where a current decision changes the quality of later states \cite{sutton2018reinforcement}. In game playing, mathematical reasoning, instruction following, and text generation, systems often generate multiple candidate paths, evaluate them, and select the option expected to lead to a better outcome. These methods are relevant to this thesis because therapeutic dialogue is also sequential: a response that appears appropriate in the current turn may still lead to resistance, disengagement, or weaker alliance in the next turn.

\subsection{Search-Based Reasoning and Lookahead Methods}

Search-based methods have been successful in domains where decisions require planning over future possibilities. In game playing, AlphaGo combined deep neural networks with tree search to evaluate possible future moves before selecting an action, showing that learned evaluation and search can support strong sequential decision-making \cite{silver2016mastering}. In language-model reasoning, Self-Consistency improves chain-of-thought reasoning by sampling multiple reasoning paths and selecting the most consistent answer rather than relying on a single greedy path \cite{wang2022selfconsistency}. Tree of Thoughts further extends this idea by allowing language models to explore, evaluate, and backtrack across multiple intermediate reasoning paths \cite{yao2023tree}. These works show that exploring alternatives before committing to a final output can improve performance in tasks that require planning or reasoning.

More recent work combines search with reward signals during LLM inference or self-training. Training Verifiers for math word problems generates multiple candidate solutions and selects the one judged most likely to be correct by a verifier model \cite{cobbe2021training}. ReST-MCTS* uses process-reward-guided Monte Carlo Tree Search to collect higher-quality reasoning traces and improve policy and reward models through self-training \cite{zhang2024restmcts}. Instead of relying only on final-answer correctness, these approaches evaluate intermediate steps or candidate paths. This suggests that future-aware selection is useful when the system must compare several possible paths before committing to one output.

\begin{figure}[htbp]
    \centering
    \begin{tikzpicture}[
        node distance=1.2cm and 1.2cm,
        every node/.style={font=\small, align=center},
        box/.style={
            rectangle,
            rounded corners,
            draw=black,
            minimum width=3.4cm,
            minimum height=0.9cm,
            text width=3.4cm,
            inner sep=5pt
        },
        smallbox/.style={
            rectangle,
            rounded corners,
            draw=black,
            minimum width=2.8cm,
            minimum height=0.8cm,
            text width=2.8cm,
            inner sep=4pt
        },
        arrow/.style={-{Latex[length=2.5mm]}, thick}
    ]

    \node[box] (state) {Current state};

    \node[smallbox, below left=1.2cm and 2.0cm of state] (c1) {Candidate 1};
    \node[smallbox, below=1.2cm of state] (c2) {Candidate 2};
    \node[smallbox, below right=1.2cm and 2.0cm of state] (c3) {Candidate 3};

    \node[smallbox, below=of c1] (f1) {Future state};
    \node[smallbox, below=of c2] (f2) {Future state};
    \node[smallbox, below=of c3] (f3) {Future state};

    \node[smallbox, below=of f1] (r1) {Reward score};
    \node[smallbox, below=of f2] (r2) {Reward score};
    \node[smallbox, below=of f3] (r3) {Reward score};

    \node[box, below=1.2cm of r2] (select) {Select highest-value path};

    \draw[arrow] (state) -- (c1);
    \draw[arrow] (state) -- (c2);
    \draw[arrow] (state) -- (c3);

    \draw[arrow] (c1) -- (f1);
    \draw[arrow] (c2) -- (f2);
    \draw[arrow] (c3) -- (f3);

    \draw[arrow] (f1) -- (r1);
    \draw[arrow] (f2) -- (r2);
    \draw[arrow] (f3) -- (r3);

    \draw[arrow] (r1) -- (select);
    \draw[arrow] (r2) -- (select);
    \draw[arrow] (r3) -- (select);

    \end{tikzpicture}
    \caption{Generic reward-guided lookahead search for response selection.}
    \label{fig:generic-lookahead-search}
\end{figure}

\subsection{Reward Models and Reward-Guided Selection}

Reward models provide a way to convert human preferences, task success, or quality judgments into a score that can guide selection or training. In reinforcement learning from human preferences, Christiano et al. trained reward models from human comparisons of trajectory segments and used them to guide agents in Atari and simulated robotics tasks \cite{christiano2017deep}. In LLM alignment, InstructGPT trained a reward model from human rankings of model outputs and used reinforcement learning from human feedback to improve instruction-following behavior \cite{ouyang2022training}. Direct Preference Optimization later showed that preference learning can also be optimized more directly without a separate reinforcement learning stage \cite{rafailov2023direct}. These works show that human or task-based preference signals can guide model behavior beyond simple next-token prediction.

Reward-guided selection has also been used directly for choosing among generated outputs. LLM-Blender uses PairRanker to compare candidate responses and GenFuser to combine high-ranked outputs, showing that ranking and selection can improve generation quality across multiple LLMs \cite{jiang2023llmblender}. In mathematical reasoning, process supervision trains reward models to evaluate intermediate reasoning steps rather than only final answers, and has been shown to outperform outcome-only supervision on challenging reasoning tasks \cite{lightman2023lets}. These methods support the idea that response quality can be improved by evaluating candidate outputs or intermediate paths before producing a final answer.

Following reward-model and ranking-based response selection approaches, candidate outputs can be scored and the highest-valued response can be selected \cite{ouyang2022training,cobbe2021training,jiang2023llmblender}:

\begin{equation}
r_t^* = \arg\max_{r_t^{(i)} \in \mathcal{R}_t} \hat{R}(H_t, r_t^{(i)})
\label{eq:generic-reward-selection}
\end{equation}

where $H_t$ is the current context, $\mathcal{R}_t$ is the set of candidate responses, $r_t^{(i)}$ is the $i$-th candidate response, $\hat{R}$ is the estimated reward function, and $r_t^*$ is the selected response.

When future outcomes are considered, the value of a candidate trajectory can be represented using the standard discounted return formulation from reinforcement learning \cite{sutton2018reinforcement}:

\begin{equation}
G_t^{(i)} = \sum_{k=0}^{d-1} \gamma^k \hat{R}_{t+k}^{(i)}
\label{eq:generic-cumulative-reward}
\end{equation}

where $G_t^{(i)}$ is the estimated value of candidate trajectory $i$, $d$ is the lookahead depth, $\gamma$ is the discount factor, and $\hat{R}_{t+k}^{(i)}$ is the predicted reward at future step $k$.

\begin{table}[htbp]
    \caption{Search and reward-guided methods relevant to future-aware response selection.}
    \label{tab:search-reward-methods}
    \small
    \begin{tabularx}{\textwidth}{p{0.18\textwidth}p{0.20\textwidth}X X}
        \toprule
        \textbf{Study} & \textbf{Area} & \textbf{Method} & \textbf{Reported Contribution} \\
        \midrule
        AlphaGo \cite{silver2016mastering} 
        & Game playing 
        & Combines deep neural networks with tree search. 
        & Shows that learned evaluation plus search can support strong sequential decision-making. \\

        Self-Consistency \cite{wang2022selfconsistency} 
        & LLM reasoning 
        & Samples multiple reasoning paths and selects the most consistent answer. 
        & Improves reasoning by avoiding reliance on a single greedy path. \\

        Tree of Thoughts \cite{yao2023tree} 
        & LLM reasoning 
        & Explores and evaluates multiple intermediate reasoning paths. 
        & Improves performance on tasks requiring planning, search, and backtracking. \\

        Training Verifiers \cite{cobbe2021training} 
        & Mathematical reasoning 
        & Generates many solutions and selects the highest-ranked one using a verifier. 
        & Shows that candidate generation plus learned verification can improve reasoning accuracy. \\

        ReST-MCTS* \cite{zhang2024restmcts} 
        & LLM self-training 
        & Uses process reward guidance with Monte Carlo Tree Search. 
        & Produces higher-quality reasoning traces and improves models through iterative self-training. \\

        InstructGPT \cite{ouyang2022training} 
        & LLM alignment 
        & Trains a reward model from human output rankings. 
        & Improves instruction-following and alignment with human preferences. \\

        LLM-Blender \cite{jiang2023llmblender} 
        & Candidate response ranking 
        & Ranks and fuses outputs from multiple LLMs. 
        & Shows that candidate ranking can improve generation quality over individual models. \\

        Process supervision \cite{lightman2023lets} 
        & Mathematical reasoning 
        & Trains reward models on intermediate reasoning steps. 
        & Shows that step-level feedback can improve reasoning reliability compared with final-answer feedback alone. \\
        \bottomrule
    \end{tabularx}
\end{table}

\subsection{Research Gap and Positioning of AIM-RGL}

The reviewed methods show that search, lookahead, ranking, and reward modeling can improve decision-making in several domains. In games, search helps evaluate future consequences before selecting a move. In reasoning tasks, lookahead allows models to compare multiple paths instead of committing to the first generated solution. In alignment and response generation, reward models and rankers help select outputs that better match human preferences or task goals.

However, these methods have mostly been applied to domains such as game playing, mathematical reasoning, instruction following, and general response generation. In LLM-based counseling, prior systems often focus on generating an appropriate current response, adapting from memory, or using observed dialogue signals. Less attention has been given to explicitly comparing possible near-future client reactions before selecting a therapeutic response. This gap is important because MI-CBT dialogue depends heavily on timing: a response may appear supportive in the current turn but still increase resistance or weaken engagement in the next turn. Therefore, this thesis positions AIM-RGL as an application of future-aware, reward-guided response selection to simulated MI-CBT therapeutic dialogue, with the detailed framework and implementation described in Chapter~3.
